% Fuente: Control 1, Pregunta 2
La cooperativa agrícola \textit{AgroCentral} le encomienda la planificación de su producción semanal con el objetivo de  maximizar el rendimiento total en toneladas. Para ello, dispone de tres cultivos: trigo, maíz y tomate, cada uno organizado en \textit{lotes}, correspondientes a parcelas estándar de la cooperativa. Considere que el trigo y el maíz poseen un canal de riego compartido, mientras que el tomate
cuenta con riego por goteo propio, cuya producción está limitada por la capacidad del invernadero. Adicionalmente, considere las siguientes características para cada
cultivo:

\medskip
\begin{center}
\renewcommand{\arraystretch}{1.3}
\begin{tabular}{lcc}
  \toprule
  \textbf{Cultivo} & \textbf{Rendimiento (ton/lote)} &
  \textbf{Turnos de riego (turnos/lote)} \\
  \midrule
  Trigo \;($x_1$)  & 1 & 2  \\
  Maíz  \;($x_2$)  & 2 & 1  \\
  Tomate\;($x_3$)  & 2 & -- \\
  \bottomrule
\end{tabular}
\end{center}

\medskip
Las condiciones operacionales son:
\begin{itemize}[leftmargin=1.5em]
  \item El canal de riego compartido dispone de un máximo de 6 turnos semanales.
  \item El invernadero tiene capacidad para un máximo de 10 lotes de tomate.
  \item Las variables de decisión pueden tomar cualquier valor real no negativo.
\end{itemize}

Por lo tanto, el problema de optimización resultante es:
\begin{equation}
  \max_{x_1,x_2,x_3}\quad x_1 + 2x_2 + 2x_3
  \qquad \text{s.a.}\quad
  2x_1 + x_2 \le 6,\quad
  x_3 \le 10,\quad
  x_1,x_2,x_3 \ge 0.
  \label{ex:4:3:eq:primal}
\end{equation}

\begin{enumerate}[label=\arabic*.]
  \item Grafique la región factible identificando todas las soluciones básicas, indicando cuáles son factibles no degeneradas, factibles degeneradas e infactibles. \textit{Hint}: El poliedro vive en $\mathbb{R}^3$.

  \item Obtenga la solución óptima usando el método \textit{Simplex Revisado}, iniciando desde $(x_1,x_2,x_3)=(0,0,0)$. Aplique la regla de Bland para elegir la variable entrante. Indique la posición de la solución dentro del conjunto factible después de cada iteración.

  \item Formule el problema dual y encuentre su solución óptima mediante el Teorema de Holgura Complementaria. Indique las unidades de medida de cada variable dual y responda: Si un agricultor vecino ofrece ceder 1 turno adicional de riego a la cooperativa. ¿Cuánto estaría dispuesto a pagar por ese turno?

  \item Actualmente el trigo rinde $c_1 = 1$ [ton/lote]. Un nuevo fertilizante podría mejorar ese rendimiento. ¿En qué rango puede variar $c_1$ sin que cambie la solución óptima actual? ¿Existe algún valor de $c_1$ para el cual convenga producir trigo?.
\end{enumerate}
